\section{Regular-prime Selmer sections, after McCallum}
\label{sec:regular-prime-sections}

The regular-prime local cusp condition combines with arithmetic-source
localization to classify the complete primary section. Throughout
this section \(p>5\) is regular and \(n=p^r\), with \(r\ge1\).
We use the all-plus model; changing the sign of \(Z\) identifies it
with the coordinates in Section~\ref{sec:irreg-setup}. Put
\(N_n=\pi_1(U_{n,\overline\Q})^{\mathrm{pro}\text{-}p}\).

\begin{lemma}[Compact cancellation and boundary centers]
\label{lem:all-regular-compact-boundary}
Let $s$ be a quotient Selmer section retaining $N_n^{\mathrm{ab}}$.
There is a unique rational coordinate cusp $P$ such that every
full local witness at $p$ is a cusp witness at $P$ and the global
compact abelian class equals that of $P$. At every finite $q\ne p$,
each local cusp witness has center $P$, and each ordinary local
witness reducing to the boundary reduces to $P$.
\end{lemma}
\begin{proof}
The power map $X_n\to X_p$ sends the compact abelian class to a
locally Kummer class. Corollary~\ref{cor:all-regular-primary-local} forces its local center at
$p$ to be a rational coordinate cusp. The fibre upstairs is boundary,
and $\Q_p$ has no nontrivial $p$-power roots of unity. Its only
$\Q_p$-point is the corresponding rational cusp $P$ of $X_n$.
An ordinary point cannot map into this boundary, so the full local
witness is cuspidal. Relative to the standard tangent at $P$, its
compact class $a\in H^1(\Q,T_pJ_n)$ is zero at $p$ and unramified
elsewhere by proper good reduction.
Proposition~\ref{prop:regular-nonabelian-localization} gives $a=0$.
For a full section its entire compact pro-$p$ image is the proper
point section $P$ by the same proposition; a smaller quotient
has only its own compact point image asserted.

For any geometric coordinate cusps $Q,P$, tangent-line ratios
have divisor $n(Q-P)$. Thus $[Q-P]$ is $p$-primary torsion. If it
were zero for distinct cusps, a function with divisor $Q-P$ would
give a degree-one map from the positive-genus curve to $\mathbb P^1$.
Hence distinct cusps have distinct classes. This calculation agrees
with the exact cusp-order statement of
\citet[p.~191, before Theorem~A]{Coleman1986}.
The $p$-primary torsion of $J_n(\Q_q)$ injects into its full
$p$-primary completion at every finite $q$: its torsion is finite
and its local group has an open additive lattice
\citep[I.3.3, p.~41]{Milne2006ADT}. Injection at an individual
finite Kummer level is not required. This proves uniqueness of
$P$ and forces any local cusp center $Q$ with zero compact
difference to equal $P$. If an ordinary point $x$ reduces to the
boundary at $q\ne p$, that boundary point lifts uniquely to
$Q\in D_n(\Q_q)$, since the boundary is finite \'etale and disjoint
there. The difference $[x-Q]$ is in a pro-$q$ formal group, on
which multiplication by $p$ is invertible. Its Kummer class is zero,
so again $Q=P$. This treats all locally occurring geometric
boundary components, without requiring interior witnesses to equal $P$.
\end{proof}

\begin{lemma}[The global abelian cusp parameter]
\label{lem:all-regular-abelian-packet}
The abelian Selmer section is the image of one global cusp section
$t_{P,\beta}$, with unique $\beta\in H^1(\Q,\Z_p(1))$. Its parameter
equals that of every full local cusp witness at $p$.
At $q\ne p$, the whole geometric inertia of a full local pro-$p$
witness is, up to conjugacy, $i_P(v_q(\beta)\tau_q)$, where
$\tau_q$ is the tame pro-$p$ character.
\end{lemma}
\begin{proof}
The inertia sequence is
\[
0\longrightarrow B_n\longrightarrow N_n^{\mathrm{ab}}
 \longrightarrow T_pJ_n\longrightarrow0,\qquad
B_n=\Z_p(1)[D_n(\overline\Q)]/\Z_p(1)\sum_Qe_Q.
\]
The compact class is zero relative to $P$.
The Tate module has no fixed vectors over $\Q$ or a finite
completion, since local torsion is finite. Thus the abelian class
has a unique preimage $b\in H^1(\Q,B_n)$, and boundary cohomology
injects locally too. Eliminate a second rational cusp generator
$e_{R_0}$ from the boundary relation. This gives an integral
$G_\Q$-equivariant splitting $B_n=I_P\oplus M_n$, where
$I_P=\Z_p(1)e_P$. Write $b=(\beta,b_M)$. The local cusp at $p$
gives $(b_M)_p=0$.

At $q\ne p$, an ordinary local witness reducing inside $U_n$
has zero geometric inertia. If it meets a boundary component,
its inertia is the tame character times its positive intersection
multiplicity. A local cusp has inertia given by its completed
Kummer valuation. Indeed, for a parameter $u$ in the strictly
henselian punctured normal disc, the pro-$p$ covers are
$u^{1/p^m}$; pullback gives the Kummer class of $u(x)$.
The same calculation applies to a cusp parameter.
Tame specialization and boundary inertia are
\citet[Expos\'e~XIII, Corollary~2.9, \S2.10 and Lemma~2.11,
pp.~288--290]{SGA1}; changing paths conjugates inertia.
By Lemma~\ref{lem:all-regular-compact-boundary}, every occurring
boundary center is $P$. Hence $b_M$ is zero on inertia away
from $p$. Coefficient inertia is trivial, so coboundaries do
not change this statement about the actual cocycle.
The module $M_n$ inherits the pro-$p$ arithmetic-source action;
Proposition~\ref{prop:regular-nonabelian-localization} gives $b_M=0$.

The summand $I_P$ and fixed-point vanishing make
$H^1(\Q,I_P)\to H^1(\Q,N_n^{\mathrm{ab}})$ injective, and likewise
at each finite place. Thus $\beta$ and its local parameter are
unique. This injection also identifies the exponent of the full
inertia formula with $v_q(\beta)$. If it is zero, the whole
inertia is zero. No finite-support restriction has been imposed.
\end{proof}

\begin{lemma}[Free pro-$p$ groups and positive-weight fixed points]
\label{lem:regular-group-fixed-points}
The group $N_n$ is free pro-$p$ of rank $n^2+1$, and each cusp
inertia is generated by a free basis element. Its ordinary
lower-central factors are finite free $\Z_p$-modules and form
the free graded Lie algebra. For nonzero $b\in\Z_p$,
$N_{N_n}(I_P^b)=I_P$.
If a quotient $R$ retains $N_n^{\mathrm{ab}}$ and all its ordinary
factors $L_m=\gamma_mR/\gamma_{m+1}R$ are torsion-free, then for
$t=t_{P,\beta}$ and every finite place $v$,
\[
 L_m^{G_{\Q_v}}=0,\qquad
 (R/\gamma_mR)^{t(G_{\Q_v})}=1,\qquad R^{t(G_\Q)}=1.
\]
The full group, its finite ordinary-central quotients and its
finite-derived quotients satisfy the factor hypothesis.
\end{lemma}
\begin{proof}
The punctured-surface presentation has $2g(X_n)$ handle generators,
$3n$ inertia generators and one relation. Eliminate a boundary
generator different from the chosen one. The free rank is
$2g+3n-1=n^2+1$; the presentation and comparison are
\citet[Expos\'e~XIII, Corollary~2.12, pp.~290--291;
Expos\'e~XII, Theorem~5.1]{SGA1}.
For every prime, the ordinary factors of a finite-rank free
pro-$p$ group are finite free with Witt ranks
\citep[\S3.5, p.~71, after Theorem~3.5.1]{LubotzkySegal2002}.
The natural free-Lie surjection has equal ranks, hence is an
isomorphism. Thus $(N_n^{\mathrm{ab}})^{\otimes m}\twoheadrightarrow L_m$.
Every pro-$p$ group's ordinary lower-central intersection is
trivial: a nonidentity element survives a finite $p$-group,
which is nilpotent.

A normalizer of $I_P^b=\langle a^b\rangle$ acts by a unit on this
procyclic group. Abelianization, where $a$ is primitive and $b\ne0$,
forces the unit to be one. In the free Lie algebra over $\Q_p$,
the centralizer of $A$ is $\Q_pA$: in the free associative algebra,
a homogeneous element commuting with $A$ is a multiple of $A^m$;
for $m>1$ this is not primitive for the characteristic-zero
coproduct. The degree-one class of a group centralizer is therefore
an integral multiple of $A$. Remove the corresponding power of
$a$. Every higher class then commutes with $bA$ and vanishes by
torsion-freeness. The trivial intersection proves the normalizer
assertion, also when $b$ is divisible by $p$.

Inner actions are trivial on $L_m$, so the cusp action there is
the original standard-tangent action. At $q\ne p$, Frobenius
eigenvalues on $T_pJ_n$ have absolute value $q^{1/2}$ and those
on the boundary are $q$ times roots of unity
\citep[Theorem~(1.6), p.~276]{Deligne1974}.
The tensor surjection gives absolute values $q^{w/2}$ on
$L_m\otimes\Q_p$, with $m\le w\le2m$; none is one.
Torsion-freeness proves fixed-point vanishing. At $p$ the graded
action factors through $\Gamma$, dominated by $G_{\Q_p}$.
Its fixed vectors are global fixed vectors, hence zero at any
good prime. Filtering $R/\gamma_mR$ gives its trivial fixed
subgroup, even for an infinitely ramified cusp parameter.
The trivial intersection also gives $R^{t(G_\Q)}=1$.

For finite-derived quotients,
\citet[\S3, Theorem~3.1, p.~101; definition p.~91]{Shmelkin1964}
gives free abelian ordinary factors for free solvable discrete
groups. For a finite nilpotent truncation, choose collected Malcev
coordinates. Its integer-valued binomial multiplication and
inversion formulas extend to $\Z_p$; density proves their laws.
Every map to a finite $p$-group is continuous: if $p^a$ is an
exponent of the target, expressions with all coordinates divisible
by $p^a$ map to one. The resulting compact pro-$p$ group is
therefore the full completion, with central factors obtained by
tensoring with $\Z_p$. Apply this to $S_d/\gamma_{c+1}S_d$,
where $F_0$ is the discrete free group on $n^2+1$ generators and
$S_d=F_0/F_0^{[d]}$ is free solvable. Its completion is
$(N_n/N_n^{[d]})/\gamma_{c+1}(N_n/N_n^{[d]})$, because both
classify the same finite $p$-quotients of the prescribed derived
length and nilpotency class. Taking $c\ge m$ proves the assertion.
No derived quotient of an open subgroup is assumed to embed in
a derived quotient downstairs.
\end{proof}

\begin{theorem}[Regular-prime cusp packets]
\label{thm:all-regular-prime-packets}
Let \(p>5\) be regular and let \(n=p^r\), \(r\ge1\).
For every stable quotient $R$ of $N_n$ retaining abelianization
and having torsion-free ordinary lower-central factors, its
point-or-cusp Selmer set is exactly the three disjoint rational
coordinate-cusp packets, each injectively parametrized by
$H^1(\Q,\Z_p(1))$. This includes the full group, every finite
ordinary-central quotient retaining abelianization and every
finite-derived quotient.
\end{theorem}
\begin{proof}
Let $t=t_{P,\beta}$ be the abelian comparison section of
Lemma~\ref{lem:all-regular-abelian-packet}.
At $q\ne p$ with $v_q(\beta)\ne0$, project a full local witness
to the full group $N_n$. Its inertia is a conjugate of
$I_P^{v_q(\beta)}$. Arithmetic inertia is normal, so the section
image normalizes this subgroup. Conjugate it to $I_P^{v_q(\beta)}$.
The standard action preserves $I_P$, and the normalizer lemma
makes its normalizer in $N_n\rtimes G_{\Q_q}$ precisely
$I_P\rtimes G_{\Q_q}$. Thus the local pro-$p$ witness is
cuspidal, with parameter $\beta_q$ by the local abelian injection.
An ordinary point in a boundary disc need not be cuspidal in the
all-prime group; only its pro-$p$ image is used here.
At $p$ the witness is actually a cusp at $P$ with parameter
$\beta_p$. These give local conjugacy to $t$ in every quotient $R$.
At the other $q\ne p$, both whole geometric inertias are zero.

Conjugate globally to make $s=t$ modulo $\gamma_2R$.
Suppose agreement holds modulo $\gamma_mR$, $m\ge2$.
The next difference is a cocycle $z_m$ in $L_m$, where the action
of $t$ equals the standard-tangent action.
At a place of local conjugacy, the image of a local conjugator
in $R/\gamma_mR$ is fixed by the common section and hence is one.
The conjugator lies in $\gamma_mR$, so its image in $L_m$
makes the local class of $z_m$ zero; in particular $(z_m)_p=0$.
At $q\ne p$ with nonzero cusp valuation, its inertia commutator
lies in $[\gamma_mR,R]\subset\gamma_{m+1}R$. Thus $z_m$ is zero
on inertia. At zero cusp valuation both whole inertias already
vanish. Prior global conjugations preserve this because the
original coefficient inertia action is trivial.

The graded cocycle is unramified outside $p$ and zero locally at
$p$. Proposition~\ref{prop:regular-nonabelian-localization} makes
its global class zero. Lift a correcting coboundary to $\gamma_mR$
and advance one step. For the original pair, the sets of conjugators
giving agreement modulo $\gamma_mR$ are nonempty, closed and nested
in the compact group $R$. A common element and the trivial central
intersection give agreement in $R$. Only the graded difference,
not the whole relative cusp cocycle, has been declared unramified.

The inertia summand injects on global abelian cohomology, so each
packet is injective in every retained quotient. Distinct centers
have distinct compact $p$-primary torsion classes.
Conversely any $\beta\in H^1(\Q,\Z_p(1))$ lifts to
$H^1(\Q,\widehat\Z(1))$ by setting all other prime components to zero.
The full global cusp section gives every required local witness,
including a real one. Odd-primary real cohomology is zero;
uniqueness of a chosen real witness is not asserted.
Parameters with infinite support are allowed, for example
valuations $b_{q_j}=p^j$ at distinct primes.
Surjectivity, injection and disjointness of the packets now prove
Theorem~\ref{thm:regular-primary-packets-intro}.
\end{proof}

\begin{corollary}[The classified prosolvable cohomology fibre]
\label{cor:classified-seventeen-kernel-fibre}
Let $n=17^r$, $r\ge1$, write $N_n^{\mathrm{sol}}$ for the
maximal geometric prosolvable group, and put
$K_{17}=\ker(N_n^{\mathrm{sol}}\to N_n)$.
Every prosolvable Selmer section has cuspidal image
$t_{P,\beta_{17}}$ in $N_n$.
Extend $\beta_{17}$ by zero at other primes to $\widehat\beta$,
and use its global prosolvable cusp section $\widehat t$ as origin.
Then
\[
 H^1(\Q,K_{17})\ \xrightarrow{\ \sim\ }\
 \operatorname{fib}_{1}\bigl(H^1(\Q,N_n^{\mathrm{sol}})
                    \longrightarrow H^1(\Q,N_n)\bigr),
\]
where all actions are induced by $\widehat t$.
Its point-or-cusp Selmer subset is exactly the image of
\[
 H^1\left(\Q,\prod_{\ell\ne17}\Z_\ell(1)\right)
\]
under the full cusp packet based at \(\widehat t\).
\end{corollary}
\begin{proof}
Every finite $17$-group is solvable, so the geometric projection
is surjective. Lift the conjugator identifying the projected
section with $t_{P,\beta_{17}}$. Its relative cocycle then lies
in $K_{17}$, proving surjectivity onto the fibre.
If two such cocycles are conjugate by $h\in N_n^{\mathrm{sol}}$,
its image in $N_n$ is fixed by the cusp action.
Lemma~\ref{lem:regular-group-fixed-points} makes that fixed subgroup
trivial. Hence $h\in K_{17}$ and the map is injective.
By Theorem~\ref{thm:regular-solvable-intro}, proved in the next
section, every Selmer section in this fibre is a full cusp section.
Its centre is \(P\), since distinct rational cusps have distinct
compact \(17\)-primary classes. The difference of its parameter
from \(\widehat\beta\) has zero \(17\)-component, and every such
prime-to-\(17\) parameter gives a Selmer section in this fibre.
The injectivity of the full cusp packet gives the asserted torsor
description. This proves
Corollary~\ref{cor:classified-seventeen-fibre-intro}.
It does not identify all of \(H^1(\Q,K_{17})\) with that subset.
\end{proof}
